\subsection{The output of the application, panel by panel}
Four layers show the effect of structure, tillage and mulch: the microscope view of the topsoil, the comparison of a mulched and a bare soil, the feedback loops of the whole field and the fate of the water saved by the mulch.

The sections are drawn from the pore-size distribution of the model; the white grains are sand, the blue spaces are the pores holding water and the dark spaces are the air-filled pores.

\appfig{v3_0}{0.96}{Tillage and structure of the topsoil, under the microscope: modelled thin sections of the coarse loam for several tillage systems.}

The bars compare the pore space of the three tillage systems and the change of drainage; consolidation closes the large pores and lowers the drainage.

\appfig{v3_1}{0.96}{Relative pore space after tillage, and the change in drainage at the current settings, coarse loam.}

The two columns show the soil with and without mulch and the bars give the evaporation, transpiration, irrigation, drainage and nitrogen leached of each.

\appfig{v4_0}{0.96}{Mulched against bare soil: evaporation, drainage and transpiration of the two surfaces.}

The line gives the difference between the setting of the user and bare soil in each flux.

\appfig{v4_1}{0.96}{Your setting against bare soil: the saving of each flux.}

The drawing is the picture of the loops of Chapter 9: rain and temperature, irrigation, the soil surface and its mulch, the crop, the drainage and the nitrogen, with the actual fluxes of the run (irrigation 250 mm, evaporation 172 mm, drainage 21 mm).

\appfig{v5_0}{0.96}{Feedback loops between climate, management, soil water, crop and nitrogen: the red arrows join the parts of the field that act on each other.}

Each row gives the effect of one step of one control (a $+0.5\,^\circ$C change, 10 per cent more rain, a stronger mulch) on the yield, the irrigation and the leaching.

\appfig{v5_1}{0.96}{What one step of each control does (click a row or a numbered arrow).}

Each panel is a question answered by sections and graphs of the model; for example the wetting and drying of the pores or the reach of the water table into the root zone.

\appfig{v5_2}{0.96}{Hidden questions of the vadose zone, answered for the coarse loam: sensors, compaction, the water table (questions 1 to 4).}

The last two panels connect the trigger to the pores that are emptied at that suction and the time that the soil delivers water to the surface.

\appfig{v5_3}{0.96}{Hidden questions of the vadose zone (questions 7 and 8): where on the suction scale the trigger really falls, and how long the soil keeps up with the sun.}

The figure divides the water that the mulch saves into less irrigation, more drainage, more transpiration and storage.

\appfig{n2_0}{0.96}{Fate of the water saved by mulch.}

The numbers are the fate of the saved millimetres: they are not all kept for the crop, and the panel follows the shares.

\appfig{n2_1}{0.96}{Fate of saved water: less irrigation 38 mm, more drainage 13 mm, more transpiration 17 mm.}

The five soil columns are the profile for five rainfall scales; the arrows show the evaporation and drainage and the colour bar gives the index that classifies the regime as drainage-dominated or evaporation-dominated.

\appfig{n3_0}{0.96}{Regime in the soil profile: vapour leaves upward, drainage downward, for rainfall scales from dry to wet.}

